% Môn: Vật lý
% Lớp: 12
% Chương: Chương I. Vật lí nhiệt
% Tên đề: Ôn tập Chương 1 - Lý 12 - Đề ôn 1 - Tân Bình - Đúng Sai
% Loại: Đúng/Sai
% Trường: THPT Tân Bình

\section{ÔN TẬP CHƯƠNG I. VẬT LÍ NHIỆT}
\subsection{Đề ôn 1 -- Tân Bình -- Đúng/Sai}

\begin{ex}
Trả lời đúng/sai khi nói về cấu tạo chất.
\choiceTF
{\True Các chất được cấu tạo từ các hạt riêng biệt gọi là phân tử.}
{Các nguyên tử, phân tử đứng sát nhau và giữa chúng không có khoảng cách.}
{\True Lực tương tác giữa các phân tử ở thể rắn lớn hơn lực tương tác giữa các phân tử ở thể lỏng và thể khí.}
{\True Các nguyên tử, phân tử chất lỏng dao động xung quanh các vị trí cân bằng không cố định.}
\loigiai{
\begin{itemchoice}
\item (Đúng) Các chất được cấu tạo từ các hạt riêng biệt; trong phạm vi kiến thức phổ thông, các hạt này có thể là nguyên tử hoặc phân tử.
\item (Sai) Giữa các nguyên tử, phân tử luôn có khoảng cách.
\item (Đúng) Ở thể rắn, lực tương tác giữa các phân tử lớn hơn so với ở thể lỏng và thể khí.
\item (Đúng) Các phân tử chất lỏng dao động và chuyển động quanh những vị trí cân bằng không cố định.
\end{itemchoice}
}
\end{ex}

\begin{ex}
Người ta truyền cho khối khí trong xi-lanh một nhiệt lượng $140\ \mathrm{J}$, làm khí nở ra và thực hiện công $200\ \mathrm{J}$. Trả lời đúng/sai với các nhận định sau.
\choiceTF
{\True Khối khí trong xi-lanh nhận nhiệt lượng $140\ \mathrm{J}$ và khí nở ra thực hiện công $200\ \mathrm{J}$.}
{\True Khối khí thực hiện công nên $A<0$ và có giá trị $A=-200\ \mathrm{J}$.}
{\True Biểu thức nguyên lí I nhiệt động lực học trong trường hợp này là $\Delta U=A+Q$.}
{Độ biến thiên nội năng của khí có giá trị là $-340\ \mathrm{J}$.}
\loigiai{
\begin{itemchoice}
\item (Đúng) Khí được truyền nhiệt lượng nên $Q=140\ \mathrm{J}$; đồng thời khí nở ra và thực hiện công $200\ \mathrm{J}$.
\item (Đúng) Theo quy ước $A$ là công mà khí nhận, khi khí thực hiện công thì $A<0$, do đó $A=-200\ \mathrm{J}$.
\item (Đúng) Nguyên lí I nhiệt động lực học được viết là $\Delta U=A+Q$.
\item (Sai) Ta có $\Delta U=A+Q=-200+140=-60\ \mathrm{J}$, không phải $-340\ \mathrm{J}$.
\end{itemchoice}
}
\end{ex}

\begin{ex}
Nhiệt nóng chảy riêng của chì là $0{,}25\times10^5\ \mathrm{J/kg}$, nhiệt độ nóng chảy của chì là $327\,^{\circ}\mathrm{C}$, biết nhiệt dung riêng của chì là $126\ \mathrm{J/(kg\cdot K)}$. Trả lời đúng/sai với các nhận định sau.
\choiceTF
{Miếng chì khối lượng $1\ \mathrm{kg}$ đang ở nhiệt độ $25\,^{\circ}\mathrm{C}$ được cung cấp nhiệt lượng $0{,}25\times10^5\ \mathrm{J}$ thì nhiệt độ của nó tăng lên $260\,^{\circ}\mathrm{C}$.}
{\True Cần cung cấp nhiệt lượng $0{,}25\times10^5\ \mathrm{J}$ để làm nóng chảy hoàn toàn $1\ \mathrm{kg}$ chì ở nhiệt độ nóng chảy của nó.}
{Cần cung cấp nhiệt lượng $25126\ \mathrm{J}$ để làm nóng chảy hoàn toàn $1\ \mathrm{kg}$ chì từ nhiệt độ $25\,^{\circ}\mathrm{C}$.}
{\True Biết công suất của lò nung là $1000\ \mathrm{W}$, giả sử hiệu suất của lò là $100\%$. Thời gian để làm nóng chảy hoàn toàn $1\ \mathrm{kg}$ chì từ nhiệt độ nóng chảy của nó bằng $25\ \mathrm{s}$.}
\loigiai{
\begin{itemchoice}
\item (Sai) Trước khi nóng chảy, nhiệt lượng làm chì tăng nhiệt độ là $Q=mc\Delta t$. Suy ra $\Delta t=\dfrac{25000}{126}\approx198{,}4\,^{\circ}\mathrm{C}$, nên nhiệt độ cuối khoảng $223{,}4\,^{\circ}\mathrm{C}$, không phải $260\,^{\circ}\mathrm{C}$.
\item (Đúng) Ở đúng nhiệt độ nóng chảy, nhiệt lượng cần cung cấp để làm nóng chảy hoàn toàn là $Q=m\lambda=1\cdot0{,}25\times10^5=25000\ \mathrm{J}$.
\item (Sai) Nhiệt lượng cần cung cấp gồm nhiệt lượng làm nóng chì đến nhiệt độ nóng chảy và nhiệt lượng nóng chảy: $Q=mc(327-25)+m\lambda=126\cdot302+25000=63052\ \mathrm{J}$.
\item (Đúng) Từ nhiệt độ nóng chảy, $Q=m\lambda=25000\ \mathrm{J}$. Với $P=1000\ \mathrm{W}$, $t=\dfrac{Q}{P}=25\ \mathrm{s}$.
\end{itemchoice}
}
\end{ex}

\begin{ex}
Một học sinh đun nóng $0{,}2\ \mathrm{kg}$ nước đá ở $0\,^{\circ}\mathrm{C}$ để hoàn toàn biến thành hơi nước ở $100\,^{\circ}\mathrm{C}$. Cho $\lambda=3{,}34\times10^5\ \mathrm{J/kg}$, $c=4{,}20\ \mathrm{kJ/(kg\cdot K)}$, $L=2{,}26\times10^6\ \mathrm{J/kg}$; bỏ qua mọi hao phí nhiệt. Trả lời đúng/sai với các nhận định sau.
\choiceTF
{\True Nhiệt lượng cần cung cấp để làm nóng chảy hoàn toàn $0{,}2\ \mathrm{kg}$ nước đá ở nhiệt độ nóng chảy là $66800\ \mathrm{J}$.}
{Nhiệt lượng cần cung cấp để đưa $0{,}2\ \mathrm{kg}$ nước từ $0\,^{\circ}\mathrm{C}$ lên $100\,^{\circ}\mathrm{C}$ là $8400\ \mathrm{J}$.}
{Nhiệt lượng cần cung cấp để hóa hơi hoàn toàn $0{,}2\ \mathrm{kg}$ nước ở $100\,^{\circ}\mathrm{C}$ là $42500\ \mathrm{J}$.}
{Tổng nhiệt lượng cần cung cấp để biến hoàn toàn $0{,}2\ \mathrm{kg}$ nước đá ở $0\,^{\circ}\mathrm{C}$ thành hơi nước ở $100\,^{\circ}\mathrm{C}$ là $60280\ \mathrm{J}$.}
\loigiai{
\begin{itemchoice}
\item (Đúng) $Q_1=m\lambda=0{,}2\cdot3{,}34\times10^5=66800\ \mathrm{J}$.
\item (Sai) $Q_2=mc\Delta t=0{,}2\cdot4200\cdot100=84000\ \mathrm{J}$, không phải $8400\ \mathrm{J}$.
\item (Sai) $Q_3=mL=0{,}2\cdot2{,}26\times10^6=452000\ \mathrm{J}$, không phải $42500\ \mathrm{J}$.
\item (Sai) Tổng nhiệt lượng là $Q=Q_1+Q_2+Q_3=66800+84000+452000=602800\ \mathrm{J}$, không phải $60280\ \mathrm{J}$.
\end{itemchoice}
}
\end{ex}
